% --- Archon physics formalization source begin ---
% archon:physics
% archon:covers PhyXMiniProblems/problem_phyx_mini_0561.lean
% archon:source-report reports/phyx_mini/problem_phyx_mini_0561.source.json

\chapter{Physics problem phyx\_mini\_0561}
\label{ch:PhyXMiniProblems_problem_phyx_mini_0561}

\paragraph{Problem source.}
\#\# Physical scenario

Assuming that the difference between Thomson's calculated \textbackslash{}( e/m \textbackslash{}) in his second experiment (see Figure 3-19) and the currently accepted value was due entirely to his neglecting the horizontal component of Earth's magnetic field outside the deflection plates.

\#\# Question

What value for that component does the difference imply? (Thomson's data: \textbackslash{}( B = 5.5 \textbackslash{}times 10\textasciicircum{}\{-4\} \textbackslash{}) T, \textbackslash{}( \textbackslash{}mathcal\{E\} = 1.5 \textbackslash{}times 10\textasciicircum{}4 \textbackslash{}) V/m, \textbackslash{}( x\_1 = 5 \textbackslash{}) cm, \textbackslash{}( y\_2/x\_2 = 8/110 \textbackslash{}).)

\#\# Answer choices

- A: A: \textbackslash{}( -18\textbackslash{}mu T \textbackslash{})

- B: B: \textbackslash{}( -43\textbackslash{}mu T \textbackslash{})

- C: C: \textbackslash{}( -22\textbackslash{}mu T \textbackslash{})

- D: D: \textbackslash{}( -31\textbackslash{}mu T \textbackslash{})

\#\# Figure caption (auxiliary; use the image as primary evidence)

The image shows a diagram of a particle's trajectory influenced by deflection plates. Here's a detailed description:

- **Trajectory**: The particle moves from left to right with an initial horizontal velocity denoted as \textbackslash{}( u\_x \textbackslash{}). As it passes between two deflection plates, its path is curved upwards, suggesting an upward force acting upon it.

- **Deflection Plates**: There are two parallel plates labeled "Deflection plates" situated along the trajectory. These plates pertain to the alteration of the particle’s path by generating an electric or magnetic field.

- **Velocity Components**: The particle has components of velocity along the x and y axes. The initial velocity is \textbackslash{}( u\_x \textbackslash{}), and as the particle deflects, an upward component \textbackslash{}( u\_y \textbackslash{}) is depicted.

- **Angles and Distances**:
  - An angle \textbackslash{}( \textbackslash{}theta \textbackslash{}) is shown between the direction of the deflected path and the horizontal.
  - Distances \textbackslash{}( x\_1 \textbackslash{}) and \textbackslash{}( x\_2 \textbackslash{}) are marked along the horizontal, showing intervals from particle entry to the deflection plates and from the plates to the screen, respectively.
  - \textbackslash{}( y\_1 \textbackslash{}) is the vertical distance between the initial trajectory and the point of deflection.
  - \textbackslash{}( y\_2 \textbackslash{}) is the vertical deflection of the particle by the time it reaches a vertical barrier on the right side.

- **Final Position**: The particle reaches a vertical screen or barrier on the right at a height \textbackslash{}( y\_2 \textbackslash{}) above the initial trajectory line.

Overall, the diagram quantitatively illustrates the effect of deflection plates on a particle's path in terms of velocity components and trajectory deflection.

\#\# Dataset metadata

Quantum Phenomena; Multi-Formula Reasoning, Physical Model Grounding Reasoning

\paragraph{Recorded answer/context.}
D: D: \textbackslash{}( -31\textbackslash{}mu T \textbackslash{})

\paragraph{Figure/image path.}
/root/proposal\_for\_physic/science-mango/phyx\_mini\_run/phyx\_data/test\_image/561.png

\paragraph{Formalization target.}
create a compiling Lean file with sorry bodies at `PhyXMiniProblems/problem\_phyx\_mini\_0561.lean`. The Lean declarations must preserve the physical quantities, dimensions or dimensional roles, figure labels, governing-law hypotheses, and final relation expressed by this problem.
Use Mathlib/Physlib names found through LeanExplore where available. If a physics API is missing, introduce faithful local abstractions rather than scalar placeholder aliases.

\begin{theorem}[Physics formalization target]
\label{thm:physics:phyx_mini_0561:target}
\lean{PhyXMiniProblems.ProblemPhyXMini0561.problem_phyx_mini_0561}
\uses{def:physics:phyx-mini-0561:phyxminiproblems-problemphyxmini0561-magneticfieldcomponentinmicroteslas, def:physics:phyx-mini-0561:phyxminiproblems-problemphyxmini0561-chargetomassincoulombsperkilogram, def:physics:phyx-mini-0561:phyxminiproblems-problemphyxmini0561-thomsonsecondexperiment, def:physics:phyx-mini-0561:phyxminiproblems-problemphyxmini0561-matchesproblemandfigurereadouts, def:physics:phyx-mini-0561:phyxminiproblems-problemphyxmini0561-hasacceptedelectronchargetomasscalibration, def:physics:phyx-mini-0561:phyxminiproblems-problemphyxmini0561-hasphysicalparameters, def:physics:phyx-mini-0561:phyxminiproblems-problemphyxmini0561-fieldsrealizescalarreadouts, def:physics:phyx-mini-0561:phyxminiproblems-problemphyxmini0561-satisfiescrossedfieldvelocityselection, def:physics:phyx-mini-0561:phyxminiproblems-problemphyxmini0561-satisfieselectricdeflectionbetweenplates, def:physics:phyx-mini-0561:phyxminiproblems-problemphyxmini0561-satisfiesearthfielddriftkinematics, def:physics:phyx-mini-0561:phyxminiproblems-problemphyxmini0561-satisfiesthomsonneglectcalculation, def:physics:phyx-mini-0561:phyxminiproblems-problemphyxmini0561-isclosestdisplayedchoice}
The assigned autoformalize agent should translate this physics problem into Lean declarations in the covered file, with theorem and lemma proof bodies written as `by sorry`.
\end{theorem}
\begin{proof}
This is an autoformalization task, not a proof task. Produce statements that can later be proved by the physics prover without weakening the physical model.
\end{proof}

% --- Archon declaration topology begin ---
\section{Declaration topology}
\begin{definition}[speed Dimension]
  \label{def:physics:phyx-mini-0561:phyxminiproblems-problemphyxmini0561-speeddimension}
  \lean{PhyXMiniProblems.ProblemPhyXMini0561.speedDimension}
  The physical dimension L T⁻¹ of a velocity component
\end{definition}
\begin{proof}
  This is the declared model definition of speed Dimension.
\end{proof}
\begin{definition}[acceleration Dimension]
  \label{def:physics:phyx-mini-0561:phyxminiproblems-problemphyxmini0561-accelerationdimension}
  \lean{PhyXMiniProblems.ProblemPhyXMini0561.accelerationDimension}
  The physical dimension L T⁻² of an acceleration component
\end{definition}
\begin{proof}
  This is the declared model definition of acceleration Dimension.
\end{proof}
\begin{definition}[electric Field Strength Dimension]
  \label{def:physics:phyx-mini-0561:phyxminiproblems-problemphyxmini0561-electricfieldstrengthdimension}
  \lean{PhyXMiniProblems.ProblemPhyXMini0561.electricFieldStrengthDimension}
  The physical dimension M L T⁻² C⁻¹ of electric-field strength
\end{definition}
\begin{proof}
  This is the declared model definition of electric Field Strength Dimension.
\end{proof}
\begin{definition}[magnetic Field Strength Dimension]
  \label{def:physics:phyx-mini-0561:phyxminiproblems-problemphyxmini0561-magneticfieldstrengthdimension}
  \lean{PhyXMiniProblems.ProblemPhyXMini0561.magneticFieldStrengthDimension}
  The physical dimension M T⁻¹ C⁻¹ of magnetic-field strength
\end{definition}
\begin{proof}
  This is the declared model definition of magnetic Field Strength Dimension.
\end{proof}
\begin{definition}[charge To Mass Dimension]
  \label{def:physics:phyx-mini-0561:phyxminiproblems-problemphyxmini0561-chargetomassdimension}
  \lean{PhyXMiniProblems.ProblemPhyXMini0561.chargeToMassDimension}
  The physical dimension C M⁻¹ of a charge-to-mass ratio
\end{definition}
\begin{proof}
  This is the declared model definition of charge To Mass Dimension.
\end{proof}
\begin{definition}[Length Quantity]
  \label{def:physics:phyx-mini-0561:phyxminiproblems-problemphyxmini0561-lengthquantity}
  \lean{PhyXMiniProblems.ProblemPhyXMini0561.LengthQuantity}
  A nonnegative physical length
\end{definition}
\begin{proof}
  This is the declared model definition of Length Quantity.
\end{proof}
\begin{definition}[Time Interval Quantity]
  \label{def:physics:phyx-mini-0561:phyxminiproblems-problemphyxmini0561-timeintervalquantity}
  \lean{PhyXMiniProblems.ProblemPhyXMini0561.TimeIntervalQuantity}
  A nonnegative physical time interval
\end{definition}
\begin{proof}
  This is the declared model definition of Time Interval Quantity.
\end{proof}
\begin{definition}[Velocity Component Quantity]
  \label{def:physics:phyx-mini-0561:phyxminiproblems-problemphyxmini0561-velocitycomponentquantity}
  \lean{PhyXMiniProblems.ProblemPhyXMini0561.VelocityComponentQuantity}
  \uses{def:physics:phyx-mini-0561:phyxminiproblems-problemphyxmini0561-speeddimension}
  A signed physical velocity component
\end{definition}
\begin{proof}
  This is the declared model definition of Velocity Component Quantity.
\end{proof}
\begin{definition}[Acceleration Component Quantity]
  \label{def:physics:phyx-mini-0561:phyxminiproblems-problemphyxmini0561-accelerationcomponentquantity}
  \lean{PhyXMiniProblems.ProblemPhyXMini0561.AccelerationComponentQuantity}
  \uses{def:physics:phyx-mini-0561:phyxminiproblems-problemphyxmini0561-accelerationdimension}
  A signed physical acceleration component
\end{definition}
\begin{proof}
  This is the declared model definition of Acceleration Component Quantity.
\end{proof}
\begin{definition}[Electric Field Strength Quantity]
  \label{def:physics:phyx-mini-0561:phyxminiproblems-problemphyxmini0561-electricfieldstrengthquantity}
  \lean{PhyXMiniProblems.ProblemPhyXMini0561.ElectricFieldStrengthQuantity}
  \uses{def:physics:phyx-mini-0561:phyxminiproblems-problemphyxmini0561-electricfieldstrengthdimension}
  A nonnegative electric-field magnitude
\end{definition}
\begin{proof}
  This is the declared model definition of Electric Field Strength Quantity.
\end{proof}
\begin{definition}[Magnetic Field Strength Quantity]
  \label{def:physics:phyx-mini-0561:phyxminiproblems-problemphyxmini0561-magneticfieldstrengthquantity}
  \lean{PhyXMiniProblems.ProblemPhyXMini0561.MagneticFieldStrengthQuantity}
  \uses{def:physics:phyx-mini-0561:phyxminiproblems-problemphyxmini0561-magneticfieldstrengthdimension}
  A nonnegative applied magnetic-field magnitude
\end{definition}
\begin{proof}
  This is the declared model definition of Magnetic Field Strength Quantity.
\end{proof}
\begin{definition}[Signed Magnetic Field Component Quantity]
  \label{def:physics:phyx-mini-0561:phyxminiproblems-problemphyxmini0561-signedmagneticfieldcomponentquantity}
  \lean{PhyXMiniProblems.ProblemPhyXMini0561.SignedMagneticFieldComponentQuantity}
  \uses{def:physics:phyx-mini-0561:phyxminiproblems-problemphyxmini0561-magneticfieldstrengthdimension}
  An oriented magnetic-field component, which may have either sign
\end{definition}
\begin{proof}
  This is the declared model definition of Signed Magnetic Field Component Quantity.
\end{proof}
\begin{definition}[Charge To Mass Magnitude Quantity]
  \label{def:physics:phyx-mini-0561:phyxminiproblems-problemphyxmini0561-chargetomassmagnitudequantity}
  \lean{PhyXMiniProblems.ProblemPhyXMini0561.ChargeToMassMagnitudeQuantity}
  \uses{def:physics:phyx-mini-0561:phyxminiproblems-problemphyxmini0561-chargetomassdimension}
  A nonnegative charge-to-mass magnitude
\end{definition}
\begin{proof}
  This is the declared model definition of Charge To Mass Magnitude Quantity.
\end{proof}
\begin{definition}[length In Meters]
  \label{def:physics:phyx-mini-0561:phyxminiproblems-problemphyxmini0561-lengthinmeters}
  \lean{PhyXMiniProblems.ProblemPhyXMini0561.lengthInMeters}
  \uses{def:physics:phyx-mini-0561:phyxminiproblems-problemphyxmini0561-lengthquantity}
  SI readout of a length, in metres
\end{definition}
\begin{proof}
  This is the declared model definition of length In Meters.
\end{proof}
\begin{definition}[time In Seconds]
  \label{def:physics:phyx-mini-0561:phyxminiproblems-problemphyxmini0561-timeinseconds}
  \lean{PhyXMiniProblems.ProblemPhyXMini0561.timeInSeconds}
  \uses{def:physics:phyx-mini-0561:phyxminiproblems-problemphyxmini0561-timeintervalquantity}
  SI readout of a time interval, in seconds
\end{definition}
\begin{proof}
  This is the declared model definition of time In Seconds.
\end{proof}
\begin{definition}[velocity In Meters Per Second]
  \label{def:physics:phyx-mini-0561:phyxminiproblems-problemphyxmini0561-velocityinmeterspersecond}
  \lean{PhyXMiniProblems.ProblemPhyXMini0561.velocityInMetersPerSecond}
  \uses{def:physics:phyx-mini-0561:phyxminiproblems-problemphyxmini0561-velocitycomponentquantity}
  SI readout of a signed velocity component, in metres per second
\end{definition}
\begin{proof}
  This is the declared model definition of velocity In Meters Per Second.
\end{proof}
\begin{definition}[acceleration In Meters Per Second Squared]
  \label{def:physics:phyx-mini-0561:phyxminiproblems-problemphyxmini0561-accelerationinmeterspersecondsquared}
  \lean{PhyXMiniProblems.ProblemPhyXMini0561.accelerationInMetersPerSecondSquared}
  \uses{def:physics:phyx-mini-0561:phyxminiproblems-problemphyxmini0561-accelerationcomponentquantity}
  SI readout of a signed acceleration, in metres per second squared
\end{definition}
\begin{proof}
  This is the declared model definition of acceleration In Meters Per Second Squared.
\end{proof}
\begin{definition}[electric Field Strength In Volts Per Meter]
  \label{def:physics:phyx-mini-0561:phyxminiproblems-problemphyxmini0561-electricfieldstrengthinvoltspermeter}
  \lean{PhyXMiniProblems.ProblemPhyXMini0561.electricFieldStrengthInVoltsPerMeter}
  \uses{def:physics:phyx-mini-0561:phyxminiproblems-problemphyxmini0561-electricfieldstrengthquantity}
  SI readout of an electric-field strength, in volts per metre
\end{definition}
\begin{proof}
  This is the declared model definition of electric Field Strength In Volts Per Meter.
\end{proof}
\begin{definition}[magnetic Field Strength In Teslas]
  \label{def:physics:phyx-mini-0561:phyxminiproblems-problemphyxmini0561-magneticfieldstrengthinteslas}
  \lean{PhyXMiniProblems.ProblemPhyXMini0561.magneticFieldStrengthInTeslas}
  \uses{def:physics:phyx-mini-0561:phyxminiproblems-problemphyxmini0561-magneticfieldstrengthquantity}
  SI readout of a magnetic-field magnitude, in teslas
\end{definition}
\begin{proof}
  This is the declared model definition of magnetic Field Strength In Teslas.
\end{proof}
\begin{definition}[magnetic Field Component In Teslas]
  \label{def:physics:phyx-mini-0561:phyxminiproblems-problemphyxmini0561-magneticfieldcomponentinteslas}
  \lean{PhyXMiniProblems.ProblemPhyXMini0561.magneticFieldComponentInTeslas}
  \uses{def:physics:phyx-mini-0561:phyxminiproblems-problemphyxmini0561-signedmagneticfieldcomponentquantity}
  SI readout of an oriented magnetic-field component, in teslas
\end{definition}
\begin{proof}
  This is the declared model definition of magnetic Field Component In Teslas.
\end{proof}
\begin{definition}[magnetic Field Component In Microteslas]
  \label{def:physics:phyx-mini-0561:phyxminiproblems-problemphyxmini0561-magneticfieldcomponentinmicroteslas}
  \lean{PhyXMiniProblems.ProblemPhyXMini0561.magneticFieldComponentInMicroteslas}
  \uses{def:physics:phyx-mini-0561:phyxminiproblems-problemphyxmini0561-signedmagneticfieldcomponentquantity, def:physics:phyx-mini-0561:phyxminiproblems-problemphyxmini0561-magneticfieldcomponentinteslas}
  Read an oriented magnetic-field component in microteslas
\end{definition}
\begin{proof}
  This is the declared model definition of magnetic Field Component In Microteslas.
\end{proof}
\begin{definition}[charge To Mass In Coulombs Per Kilogram]
  \label{def:physics:phyx-mini-0561:phyxminiproblems-problemphyxmini0561-chargetomassincoulombsperkilogram}
  \lean{PhyXMiniProblems.ProblemPhyXMini0561.chargeToMassInCoulombsPerKilogram}
  \uses{def:physics:phyx-mini-0561:phyxminiproblems-problemphyxmini0561-chargetomassmagnitudequantity}
  SI readout of a charge-to-mass magnitude, in coulombs per kilogram
\end{definition}
\begin{proof}
  This is the declared model definition of charge To Mass In Coulombs Per Kilogram.
\end{proof}
\begin{definition}[Particle Species]
  \label{def:physics:phyx-mini-0561:phyxminiproblems-problemphyxmini0561-particlespecies}
  \lean{PhyXMiniProblems.ProblemPhyXMini0561.ParticleSpecies}
  Particle-species roles distinguished in the experiment
\end{definition}
\begin{proof}
  This is the declared model definition of Particle Species.
\end{proof}
\begin{definition}[Plate Label]
  \label{def:physics:phyx-mini-0561:phyxminiproblems-problemphyxmini0561-platelabel}
  \lean{PhyXMiniProblems.ProblemPhyXMini0561.PlateLabel}
  The two deflection plates drawn above and below the beam
\end{definition}
\begin{proof}
  This is the declared model definition of Plate Label.
\end{proof}
\begin{definition}[Plate Arrangement]
  \label{def:physics:phyx-mini-0561:phyxminiproblems-problemphyxmini0561-platearrangement}
  \lean{PhyXMiniProblems.ProblemPhyXMini0561.PlateArrangement}
  Whether the two deflection plates have the pictured parallel geometry
\end{definition}
\begin{proof}
  This is the declared model definition of Plate Arrangement.
\end{proof}
\begin{definition}[Figure Axis]
  \label{def:physics:phyx-mini-0561:phyxminiproblems-problemphyxmini0561-figureaxis}
  \lean{PhyXMiniProblems.ProblemPhyXMini0561.FigureAxis}
  The two labeled axes of the trajectory diagram
\end{definition}
\begin{proof}
  This is the declared model definition of Figure Axis.
\end{proof}
\begin{definition}[Trajectory Station]
  \label{def:physics:phyx-mini-0561:phyxminiproblems-problemphyxmini0561-trajectorystation}
  \lean{PhyXMiniProblems.ProblemPhyXMini0561.TrajectoryStation}
  Named positions along the trajectory in the supplied figure
\end{definition}
\begin{proof}
  This is the declared model definition of Trajectory Station.
\end{proof}
\begin{definition}[Horizontal Beam Direction]
  \label{def:physics:phyx-mini-0561:phyxminiproblems-problemphyxmini0561-horizontalbeamdirection}
  \lean{PhyXMiniProblems.ProblemPhyXMini0561.HorizontalBeamDirection}
  The direction of the beam arrow labeled uₓ
\end{definition}
\begin{proof}
  This is the declared model definition of Horizontal Beam Direction.
\end{proof}
\begin{definition}[beam Direction]
  \label{def:physics:phyx-mini-0561:phyxminiproblems-problemphyxmini0561-beamdirection}
  \lean{PhyXMiniProblems.ProblemPhyXMini0561.beamDirection}
  The unit direction of the rightward beam in three-dimensional space
\end{definition}
\begin{proof}
  This is the declared model definition of beam Direction.
\end{proof}
\begin{definition}[upward Direction]
  \label{def:physics:phyx-mini-0561:phyxminiproblems-problemphyxmini0561-upwarddirection}
  \lean{PhyXMiniProblems.ProblemPhyXMini0561.upwardDirection}
  The unit direction drawn upward in the trajectory plane
\end{definition}
\begin{proof}
  This is the declared model definition of upward Direction.
\end{proof}
\begin{definition}[earth Horizontal Direction]
  \label{def:physics:phyx-mini-0561:phyxminiproblems-problemphyxmini0561-earthhorizontaldirection}
  \lean{PhyXMiniProblems.ProblemPhyXMini0561.earthHorizontalDirection}
  The horizontal direction normal to the vertical trajectory plane. This is the relevant component direction for Earth's magnetic field
\end{definition}
\begin{proof}
  This is the declared model definition of earth Horizontal Direction.
\end{proof}
\begin{definition}[Thomson Second Experiment]
  \label{def:physics:phyx-mini-0561:phyxminiproblems-problemphyxmini0561-thomsonsecondexperiment}
  \lean{PhyXMiniProblems.ProblemPhyXMini0561.ThomsonSecondExperiment}
  \uses{def:physics:phyx-mini-0561:phyxminiproblems-problemphyxmini0561-lengthquantity, def:physics:phyx-mini-0561:phyxminiproblems-problemphyxmini0561-timeintervalquantity, def:physics:phyx-mini-0561:phyxminiproblems-problemphyxmini0561-velocitycomponentquantity, def:physics:phyx-mini-0561:phyxminiproblems-problemphyxmini0561-accelerationcomponentquantity, def:physics:phyx-mini-0561:phyxminiproblems-problemphyxmini0561-electricfieldstrengthquantity, def:physics:phyx-mini-0561:phyxminiproblems-problemphyxmini0561-magneticfieldstrengthquantity, def:physics:phyx-mini-0561:phyxminiproblems-problemphyxmini0561-signedmagneticfieldcomponentquantity, def:physics:phyx-mini-0561:phyxminiproblems-problemphyxmini0561-chargetomassmagnitudequantity, def:physics:phyx-mini-0561:phyxminiproblems-problemphyxmini0561-particlespecies, def:physics:phyx-mini-0561:phyxminiproblems-problemphyxmini0561-platelabel, def:physics:phyx-mini-0561:phyxminiproblems-problemphyxmini0561-platearrangement, def:physics:phyx-mini-0561:phyxminiproblems-problemphyxmini0561-figureaxis, def:physics:phyx-mini-0561:phyxminiproblems-problemphyxmini0561-trajectorystation, def:physics:phyx-mini-0561:phyxminiproblems-problemphyxmini0561-horizontalbeamdirection}
  Independent physical quantities and labeled geometric observables in Thomson's experiment. In particular, earthHorizontalComponent is an unknown observable: its requested value is not stored elsewhere in the setup
\end{definition}
\begin{proof}
  This is the declared model definition of Thomson Second Experiment.
\end{proof}
\begin{definition}[Matches Problem And Figure Readouts]
  \label{def:physics:phyx-mini-0561:phyxminiproblems-problemphyxmini0561-matchesproblemandfigurereadouts}
  \lean{PhyXMiniProblems.ProblemPhyXMini0561.MatchesProblemAndFigureReadouts}
  \uses{def:physics:phyx-mini-0561:phyxminiproblems-problemphyxmini0561-lengthinmeters, def:physics:phyx-mini-0561:phyxminiproblems-problemphyxmini0561-velocityinmeterspersecond, def:physics:phyx-mini-0561:phyxminiproblems-problemphyxmini0561-electricfieldstrengthinvoltspermeter, def:physics:phyx-mini-0561:phyxminiproblems-problemphyxmini0561-magneticfieldstrengthinteslas, def:physics:phyx-mini-0561:phyxminiproblems-problemphyxmini0561-thomsonsecondexperiment}
  Problem data and qualitative information read directly from the primary figure. The displayed y₂ / x₂ = 8 / 110 is represented by the conventional centimetre readouts y₂ = 8 cm and x₂ = 110 cm, together with the quotient. The image shows that y₂ is measured from the original beam line, not from the plate-exit point at height y₁. This predicate contains no value for Earth's field and no answer-choice assertion
\end{definition}
\begin{proof}
  This is the declared model definition of Matches Problem And Figure Readouts.
\end{proof}
\begin{definition}[Has Accepted Electron Charge To Mass Calibration]
  \label{def:physics:phyx-mini-0561:phyxminiproblems-problemphyxmini0561-hasacceptedelectronchargetomasscalibration}
  \lean{PhyXMiniProblems.ProblemPhyXMini0561.HasAcceptedElectronChargeToMassCalibration}
  \uses{def:physics:phyx-mini-0561:phyxminiproblems-problemphyxmini0561-chargetomassincoulombsperkilogram, def:physics:phyx-mini-0561:phyxminiproblems-problemphyxmini0561-thomsonsecondexperiment}
  The modern accepted magnitude of the electron charge-to-mass ratio used to interpret the historical discrepancy. This is external calibration data, not a statement about the requested terrestrial magnetic field
\end{definition}
\begin{proof}
  This is the declared model definition of Has Accepted Electron Charge To Mass Calibration.
\end{proof}
\begin{definition}[Has Physical Parameters]
  \label{def:physics:phyx-mini-0561:phyxminiproblems-problemphyxmini0561-hasphysicalparameters}
  \lean{PhyXMiniProblems.ProblemPhyXMini0561.HasPhysicalParameters}
  \uses{def:physics:phyx-mini-0561:phyxminiproblems-problemphyxmini0561-lengthinmeters, def:physics:phyx-mini-0561:phyxminiproblems-problemphyxmini0561-timeinseconds, def:physics:phyx-mini-0561:phyxminiproblems-problemphyxmini0561-velocityinmeterspersecond, def:physics:phyx-mini-0561:phyxminiproblems-problemphyxmini0561-electricfieldstrengthinvoltspermeter, def:physics:phyx-mini-0561:phyxminiproblems-problemphyxmini0561-magneticfieldstrengthinteslas, def:physics:phyx-mini-0561:phyxminiproblems-problemphyxmini0561-chargetomassincoulombsperkilogram, def:physics:phyx-mini-0561:phyxminiproblems-problemphyxmini0561-thomsonsecondexperiment}
  Positivity and nondegeneracy conditions for the physical experiment
\end{definition}
\begin{proof}
  This is the declared model definition of Has Physical Parameters.
\end{proof}
\begin{definition}[Field Vanishes Outside]
  \label{def:physics:phyx-mini-0561:phyxminiproblems-problemphyxmini0561-fieldvanishesoutside}
  \lean{PhyXMiniProblems.ProblemPhyXMini0561.FieldVanishesOutside}
  A spacetime-dependent vector field vanishes outside a specified region
\end{definition}
\begin{proof}
  This is the declared model definition of Field Vanishes Outside.
\end{proof}
\begin{definition}[Fields Realize Scalar Readouts]
  \label{def:physics:phyx-mini-0561:phyxminiproblems-problemphyxmini0561-fieldsrealizescalarreadouts}
  \lean{PhyXMiniProblems.ProblemPhyXMini0561.FieldsRealizeScalarReadouts}
  \uses{def:physics:phyx-mini-0561:phyxminiproblems-problemphyxmini0561-electricfieldstrengthinvoltspermeter, def:physics:phyx-mini-0561:phyxminiproblems-problemphyxmini0561-magneticfieldstrengthinteslas, def:physics:phyx-mini-0561:phyxminiproblems-problemphyxmini0561-magneticfieldcomponentinteslas, def:physics:phyx-mini-0561:phyxminiproblems-problemphyxmini0561-upwarddirection, def:physics:phyx-mini-0561:phyxminiproblems-problemphyxmini0561-earthhorizontaldirection, def:physics:phyx-mini-0561:phyxminiproblems-problemphyxmini0561-thomsonsecondexperiment, def:physics:phyx-mini-0561:phyxminiproblems-problemphyxmini0561-fieldvanishesoutside}
  The Physlib vector fields realize the stored scalar SI readouts. Applied selector fields are confined to the plates. The terrestrial field is modeled in the outside drift region with its independent signed coefficient along the horizontal direction normal to the trajectory plane. The applied selector field points opposite that positive direction, so its magnetic force on a right-moving electron opposes the upward electric force
\end{definition}
\begin{proof}
  This is the declared model definition of Fields Realize Scalar Readouts.
\end{proof}
\begin{definition}[Satisfies Crossed Field Velocity Selection]
  \label{def:physics:phyx-mini-0561:phyxminiproblems-problemphyxmini0561-satisfiescrossedfieldvelocityselection}
  \lean{PhyXMiniProblems.ProblemPhyXMini0561.SatisfiesCrossedFieldVelocitySelection}
  \uses{def:physics:phyx-mini-0561:phyxminiproblems-problemphyxmini0561-velocityinmeterspersecond, def:physics:phyx-mini-0561:phyxminiproblems-problemphyxmini0561-electricfieldstrengthinvoltspermeter, def:physics:phyx-mini-0561:phyxminiproblems-problemphyxmini0561-magneticfieldstrengthinteslas, def:physics:phyx-mini-0561:phyxminiproblems-problemphyxmini0561-thomsonsecondexperiment}
  Crossed-field velocity selection: equality of electric and magnetic Lorentz force magnitudes gives E = uₓ B. This law contains no terrestrial-field value
\end{definition}
\begin{proof}
  This is the declared model definition of Satisfies Crossed Field Velocity Selection.
\end{proof}
\begin{definition}[Satisfies Electric Deflection Between Plates]
  \label{def:physics:phyx-mini-0561:phyxminiproblems-problemphyxmini0561-satisfieselectricdeflectionbetweenplates}
  \lean{PhyXMiniProblems.ProblemPhyXMini0561.SatisfiesElectricDeflectionBetweenPlates}
  \uses{def:physics:phyx-mini-0561:phyxminiproblems-problemphyxmini0561-lengthinmeters, def:physics:phyx-mini-0561:phyxminiproblems-problemphyxmini0561-timeinseconds, def:physics:phyx-mini-0561:phyxminiproblems-problemphyxmini0561-velocityinmeterspersecond, def:physics:phyx-mini-0561:phyxminiproblems-problemphyxmini0561-accelerationinmeterspersecondsquared, def:physics:phyx-mini-0561:phyxminiproblems-problemphyxmini0561-electricfieldstrengthinvoltspermeter, def:physics:phyx-mini-0561:phyxminiproblems-problemphyxmini0561-chargetomassincoulombsperkilogram, def:physics:phyx-mini-0561:phyxminiproblems-problemphyxmini0561-thomsonsecondexperiment}
  Uniform horizontal motion and constant upward electric acceleration between the plates. The acceleration is the electric Lorentz-force magnitude divided by mass, a\_E = (|e|/m) E
\end{definition}
\begin{proof}
  This is the declared model definition of Satisfies Electric Deflection Between Plates.
\end{proof}
\begin{definition}[Satisfies Earth Field Drift Kinematics]
  \label{def:physics:phyx-mini-0561:phyxminiproblems-problemphyxmini0561-satisfiesearthfielddriftkinematics}
  \lean{PhyXMiniProblems.ProblemPhyXMini0561.SatisfiesEarthFieldDriftKinematics}
  \uses{def:physics:phyx-mini-0561:phyxminiproblems-problemphyxmini0561-lengthinmeters, def:physics:phyx-mini-0561:phyxminiproblems-problemphyxmini0561-timeinseconds, def:physics:phyx-mini-0561:phyxminiproblems-problemphyxmini0561-velocityinmeterspersecond, def:physics:phyx-mini-0561:phyxminiproblems-problemphyxmini0561-accelerationinmeterspersecondsquared, def:physics:phyx-mini-0561:phyxminiproblems-problemphyxmini0561-magneticfieldcomponentinteslas, def:physics:phyx-mini-0561:phyxminiproblems-problemphyxmini0561-chargetomassincoulombsperkilogram, def:physics:phyx-mini-0561:phyxminiproblems-problemphyxmini0561-thomsonsecondexperiment}
  Outside the plates, Earth's horizontal component is the only additional source of vertical acceleration. Its signed Lorentz acceleration is a\_B = (|e|/m) uₓ B\_E. The figure's y₂ is the total height above the original beam line, so the drift law begins from the nonzero exit height y₁; it is not a displacement measured from the plate exit
\end{definition}
\begin{proof}
  This is the declared model definition of Satisfies Earth Field Drift Kinematics.
\end{proof}
\begin{definition}[Satisfies Thomson Neglect Calculation]
  \label{def:physics:phyx-mini-0561:phyxminiproblems-problemphyxmini0561-satisfiesthomsonneglectcalculation}
  \lean{PhyXMiniProblems.ProblemPhyXMini0561.SatisfiesThomsonNeglectCalculation}
  \uses{def:physics:phyx-mini-0561:phyxminiproblems-problemphyxmini0561-lengthinmeters, def:physics:phyx-mini-0561:phyxminiproblems-problemphyxmini0561-velocityinmeterspersecond, def:physics:phyx-mini-0561:phyxminiproblems-problemphyxmini0561-electricfieldstrengthinvoltspermeter, def:physics:phyx-mini-0561:phyxminiproblems-problemphyxmini0561-chargetomassincoulombsperkilogram, def:physics:phyx-mini-0561:phyxminiproblems-problemphyxmini0561-thomsonsecondexperiment}
  Thomson's neglect calculation sets the outside magnetic acceleration to zero. Eliminating y₁, uᵧ, and the transit times from the electric-only model gives (|e|/m)\_T = y₂ uₓ² / (E x₁ (x₂ + x₁/2)). This historical calculation law gives neither a numerical discrepancy nor a value for Earth's field
\end{definition}
\begin{proof}
  This is the declared model definition of Satisfies Thomson Neglect Calculation.
\end{proof}
\begin{definition}[Answer Choice]
  \label{def:physics:phyx-mini-0561:phyxminiproblems-problemphyxmini0561-answerchoice}
  \lean{PhyXMiniProblems.ProblemPhyXMini0561.AnswerChoice}
  The four signed horizontal field components printed as answer choices
\end{definition}
\begin{proof}
  This is the declared model definition of Answer Choice.
\end{proof}
\begin{definition}[displayed Component In Microteslas]
  \label{def:physics:phyx-mini-0561:phyxminiproblems-problemphyxmini0561-displayedcomponentinmicroteslas}
  \lean{PhyXMiniProblems.ProblemPhyXMini0561.displayedComponentInMicroteslas}
  \uses{def:physics:phyx-mini-0561:phyxminiproblems-problemphyxmini0561-answerchoice}
  Signed magnetic-field component, in microteslas, displayed by a choice
\end{definition}
\begin{proof}
  This is the declared model definition of displayed Component In Microteslas.
\end{proof}
\begin{definition}[Is Closest Displayed Choice]
  \label{def:physics:phyx-mini-0561:phyxminiproblems-problemphyxmini0561-isclosestdisplayedchoice}
  \lean{PhyXMiniProblems.ProblemPhyXMini0561.IsClosestDisplayedChoice}
  \uses{def:physics:phyx-mini-0561:phyxminiproblems-problemphyxmini0561-magneticfieldcomponentinmicroteslas, def:physics:phyx-mini-0561:phyxminiproblems-problemphyxmini0561-thomsonsecondexperiment, def:physics:phyx-mini-0561:phyxminiproblems-problemphyxmini0561-answerchoice, def:physics:phyx-mini-0561:phyxminiproblems-problemphyxmini0561-displayedcomponentinmicroteslas}
  A choice is at least as close to the inferred component as every choice
\end{definition}
\begin{proof}
  This is the declared model definition of Is Closest Displayed Choice.
\end{proof}
\begin{lemma}[earth Horizontal Component exact Formula]
  \label{lem:physics:phyx-mini-0561:phyxminiproblems-problemphyxmini0561-earthhorizontalcomponent-exactformula}
  \lean{PhyXMiniProblems.ProblemPhyXMini0561.earthHorizontalComponent_exactFormula}
  \uses{def:physics:phyx-mini-0561:phyxminiproblems-problemphyxmini0561-lengthinmeters, def:physics:phyx-mini-0561:phyxminiproblems-problemphyxmini0561-velocityinmeterspersecond, def:physics:phyx-mini-0561:phyxminiproblems-problemphyxmini0561-electricfieldstrengthinvoltspermeter, def:physics:phyx-mini-0561:phyxminiproblems-problemphyxmini0561-magneticfieldcomponentinteslas, def:physics:phyx-mini-0561:phyxminiproblems-problemphyxmini0561-chargetomassincoulombsperkilogram, def:physics:phyx-mini-0561:phyxminiproblems-problemphyxmini0561-thomsonsecondexperiment, def:physics:phyx-mini-0561:phyxminiproblems-problemphyxmini0561-matchesproblemandfigurereadouts, def:physics:phyx-mini-0561:phyxminiproblems-problemphyxmini0561-hasphysicalparameters, def:physics:phyx-mini-0561:phyxminiproblems-problemphyxmini0561-satisfiescrossedfieldvelocityselection, def:physics:phyx-mini-0561:phyxminiproblems-problemphyxmini0561-satisfieselectricdeflectionbetweenplates, def:physics:phyx-mini-0561:phyxminiproblems-problemphyxmini0561-satisfiesearthfielddriftkinematics}
  Eliminating the two transit times and accelerations gives the general correction formula. The mean vertical slope during the outside drift is (y₂ - y₁) / x₂; the second term is the true plate-exit slope
\end{lemma}
\begin{proof}
  The conclusion follows from the governing relations and figure readouts named in the statement of earth Horizontal Component exact Formula.
\end{proof}
\begin{lemma}[thomson Neglect Charge To Mass readout]
  \label{lem:physics:phyx-mini-0561:phyxminiproblems-problemphyxmini0561-thomsonneglectchargetomass-readout}
  \lean{PhyXMiniProblems.ProblemPhyXMini0561.thomsonNeglectChargeToMass_readout}
  \uses{def:physics:phyx-mini-0561:phyxminiproblems-problemphyxmini0561-chargetomassincoulombsperkilogram, def:physics:phyx-mini-0561:phyxminiproblems-problemphyxmini0561-thomsonsecondexperiment, def:physics:phyx-mini-0561:phyxminiproblems-problemphyxmini0561-matchesproblemandfigurereadouts, def:physics:phyx-mini-0561:phyxminiproblems-problemphyxmini0561-hasphysicalparameters, def:physics:phyx-mini-0561:phyxminiproblems-problemphyxmini0561-satisfiescrossedfieldvelocityselection, def:physics:phyx-mini-0561:phyxminiproblems-problemphyxmini0561-satisfiesthomsonneglectcalculation}
  The electric-only calculation gives Thomson's inferred value, which is below the accepted calibration. This makes the discrepancy named in the question explicit without assuming the requested Earth-field correction
\end{lemma}
\begin{proof}
  The conclusion follows from the governing relations and figure readouts named in the statement of thomson Neglect Charge To Mass readout.
\end{proof}
\begin{lemma}[earth Horizontal Component microteslas]
  \label{lem:physics:phyx-mini-0561:phyxminiproblems-problemphyxmini0561-earthhorizontalcomponent-microteslas}
  \lean{PhyXMiniProblems.ProblemPhyXMini0561.earthHorizontalComponent_microteslas}
  \uses{def:physics:phyx-mini-0561:phyxminiproblems-problemphyxmini0561-magneticfieldcomponentinmicroteslas, def:physics:phyx-mini-0561:phyxminiproblems-problemphyxmini0561-thomsonsecondexperiment, def:physics:phyx-mini-0561:phyxminiproblems-problemphyxmini0561-matchesproblemandfigurereadouts, def:physics:phyx-mini-0561:phyxminiproblems-problemphyxmini0561-hasacceptedelectronchargetomasscalibration, def:physics:phyx-mini-0561:phyxminiproblems-problemphyxmini0561-hasphysicalparameters, def:physics:phyx-mini-0561:phyxminiproblems-problemphyxmini0561-satisfiescrossedfieldvelocityselection, def:physics:phyx-mini-0561:phyxminiproblems-problemphyxmini0561-satisfieselectricdeflectionbetweenplates, def:physics:phyx-mini-0561:phyxminiproblems-problemphyxmini0561-satisfiesearthfielddriftkinematics}
  With the supplied readouts and accepted electron ratio, the inferred terrestrial component is exactly -897375 / 29282 μT, approximately -30.646 μT. The numerical field value remains a conclusion, not a calibration or law premise
\end{lemma}
\begin{proof}
  The conclusion follows from the governing relations and figure readouts named in the statement of earth Horizontal Component microteslas.
\end{proof}
% --- Archon declaration topology end ---
% --- Archon physics formalization source end ---
